\section{Feature Engineering}\label{sec:method:features}

The scaled \glspl{pdw} of \cref{sec:method:data} already carry carrier frequency, pulse width, amplitude, absolute \gls{toa}, and \gls{aoa}, but they do not make the local pulse-repetition structure explicit.
Emitters in agile or staggered modes are often easier to tell apart once inter-arrival statistics are exposed, because stable or patterned \gls{pri} schedules leave distinct signatures in successive $\Delta\mathrm{TOA}$ values even when absolute timing has been min--max mapped.
This section therefore defines a small set of causal, hand-crafted timing features, documents how they are appended to the base input, and states the burn-in rule that makes every engineered channel well defined.

\subsection{Delta \gls{toa}}\label{sec:method:features:delta}

Let $t_n$ denote the physical \gls{toa} of pulse $n$ on a recording, taken from the simulator export before the min--max map of \cref{eq:method:minmax}.
Successive inter-arrival intervals follow the same convention as in \cref{eq:bg:pdw_pri},
\begin{equation}\label{eq:method:delta_toa}
  \tau_n = t_n - t_{n-1},
  \qquad n \ge 2,
\end{equation}
so $\tau_n$ is the instantaneous $\Delta\mathrm{TOA}$ (equivalently the local \gls{pri}) between pulses $n-1$ and $n$.
Pulse $n=1$ has no predecessor and is excluded from all quantities below that depend on $\tau_n$.

\subsection{Rolling summaries}\label{sec:method:features:rolling}

Agile and staggered waveforms modulate $\tau_n$ over short and long horizons.
To expose both fast jitter and slower schedule changes, each pulse $n \ge 2$ is assigned rolling summaries over the three look-back widths $k \in \mathcal{K} = \{4,16,64\}$,
\begin{equation}\label{eq:method:rolling_mean}
  \bar{\tau}^{(k)}_n
  =
  \frac{1}{k}\sum_{r=0}^{k-1} \tau_{n-r},
\end{equation}
\begin{equation}\label{eq:method:rolling_std}
  \hat{\sigma}^{(k)}_n
  =
  \sqrt{\frac{1}{k}\sum_{r=0}^{k-1}\bigl(\tau_{n-r} - \bar{\tau}^{(k)}_n\bigr)^2}.
\end{equation}
The mean $\bar{\tau}^{(k)}_n$ tracks the recent average $\Delta\mathrm{TOA}$, whereas $\hat{\sigma}^{(k)}_n$ captures local variability in the same window.
The three window sizes probe complementary time scales: $k=4$ reacts to short bursts of agility, $k=16$ summarises medium-range stagger patterns, and $k=64$ approximates longer-term repetition behaviour within a pulse train.

For a window of width $k$ to contain $k$ valid $\Delta\mathrm{TOA}$ values, pulse index $n$ must satisfy $n \ge k + 1$ (because $\tau_1$ is undefined).
The widest window therefore requires $n \ge 65$.

\subsection{Burn-in and augmented input}\label{sec:method:features:burnin}

When feature engineering is enabled, the first $B = 64$ pulses of every recording are discarded before windowing (\cref{sec:method:data:window}).
This burn-in guarantees that every retained pulse satisfies $n \ge 65$ and therefore carries a full set of rolling statistics for all $k \in \mathcal{K}$, together with its instantaneous $\tau_n$.
Pulses removed by burn-in are not used for training, validation, or test metrics in the engineered-feature configuration; the baseline configuration that omits engineered channels keeps all pulses and applies windowing directly to the scaled base \glspl{pdw}.

Write $\tilde{\vect{x}}^{\mathrm{base}}_n \in \R^{d_{\mathrm{base}}}$ for the per-pulse vector after the scaling rules of \cref{sec:method:data:norm} ($d_{\mathrm{base}} = 6$ for carrier frequency, pulse width, min--max \gls{toa}, $\log_{10}$ amplitude, and the two rescaled \gls{aoa} components).
The engineered timing block is
\begin{equation}\label{eq:method:engineered_block}
  \vect{f}_n
  =
  \bigl(
    \tau_n,\,
    \bar{\tau}^{(4)}_n,\,
    \bar{\tau}^{(16)}_n,\,
    \bar{\tau}^{(64)}_n,\,
    \hat{\sigma}^{(4)}_n,\,
    \hat{\sigma}^{(16)}_n,\,
    \hat{\sigma}^{(64)}_n
  \bigr)^{\top}
  \in \R^{d_{\mathrm{fe}}},
  \qquad
  d_{\mathrm{fe}} = 7,
\end{equation}
and the augmented pulse vector fed to the encoder is
\begin{equation}\label{eq:method:augmented_input}
  \vect{x}_n
  =
  \begin{cases}
    \tilde{\vect{x}}^{\mathrm{base}}_n,
    & \text{without feature engineering}, \\[4pt]
    \bigl(\tilde{\vect{x}}^{\mathrm{base}}_n,\, \tilde{\vect{f}}_n\bigr),
    & \text{with feature engineering},
  \end{cases}
\end{equation}
where $\tilde{\vect{f}}_n$ applies the same training-corpus $z$-score standardisation of \cref{eq:method:standardize,eq:method:standardize_stats} to each coordinate of $\vect{f}_n$, with means and variances estimated on all retained training pulses ($n \ge 65$ after burn-in) and applied unchanged to validation and test pulses.
The input dimension is therefore $d = d_{\mathrm{base}}$ in the baseline setting and $d = d_{\mathrm{base}} + d_{\mathrm{fe}} = 13$ when engineered timing features are enabled; the affine input embedding of \cref{eq:method:input_embed} is instantiated with the matching column count in each case.

\subsection{Experimental comparison}\label{sec:method:features:ablation}

Feature engineering is treated as an optional preprocessing stage rather than a change to the encoder topology.
Every other architectural and training choice in \cref{sec:method:arch,sec:method:loss,sec:method:training} is held fixed, and two variants are trained and evaluated on the same scenario corpus:
\begin{itemize}
  \item \textbf{Baseline inputs} ($d = d_{\mathrm{base}}$): scaled \glspl{pdw} only, no burn-in.
  \item \textbf{Engineered inputs} ($d = d_{\mathrm{base}} + d_{\mathrm{fe}}$): concatenation of scaled \glspl{pdw} and the seven timing features of \cref{eq:method:engineered_block}, with a per-recording burn-in of $B = 64$ pulses.
\end{itemize}
We will report the paired comparison in \cref{sec:exp:ablation} to quantify how much of the emitter and mode clustering performance is already captured by the transformer over raw \glspl{pdw} versus how much is recovered by exposing local \gls{pri} structure explicitly.
